\section{引言}
\label{sec:introduction}

二维光子晶体中的线缺陷破坏了横向的完整周期结构，却保留了沿缺陷方向的周期性。
当频率落在背景周期介质的谱隙内时，这种缺陷可能支持沿缺陷传播而在横向局域的导模。
相应问题既不是完整周期介质的普通能带计算，也不是均匀开放波导的横截面特征值问题；
它是在无界周期背景中寻找缺陷诱导的离散谱。Fliss 利用准周期谱分解和
Dirichlet-to-Neumann 映射，把这类导模刻画为固定 Bloch 波数下的条带特征值问题
\cite{Fliss2013}。

开放周期端的精确处理和高精度离散是导模计算的两个核心困难。边界积分方法能够只离散
材料界面，并已被用于构造高阶波导模态求解器\cite{LuLu2014}；但一个很小的离散矩阵残量
并不自动给出连续物理特征值的可靠位数。Barnett、Hassell 与 Tacy 以边界量控制谱距离的
工作说明，边界上的不相容量可以成为谱包含的有效入口\cite{BarnettHassellTacy2018}；不过
其有界光滑 Neumann 问题与本文的无界周期传输问题不同，不能直接给出本文的结论。光子晶体
特征值问题中的有限元研究也表明，后验估计可以用于衡量特征值误差并指导离散加密
\cite{EngstromGianiGrubisic2016}。本文为当前条带算子另外证明一个预解式谱距离定理，并用
经过向外舍入的区间算术给出各项输入的可靠单侧界。

本文研究两侧为周期介质、中心为一般局部缺陷的二维波导。每个含材料的基本胞元至多含一个
单连通介质柱截面；中心胞元可以含一个形状不同的介质柱，也可以为空。后者只是本文数值
算例采用的特例，而不是一般模型的限制。与把材料写成连续折射率函数的通常表述相比，本文
采用具有尖锐界面的分片常数折射率，因此偏微分方程必须在各均匀子区域内理解，并在每条
材料界面上显式满足传输条件。

本文首先在固定 Bloch 波数下用边界积分、端口 Fourier 截断和有序 QZ 分解计算候选
$\widehat k$ 及其分片重构场。随后固定同一候选、密度、传播矩阵、材料和几何输入，直接
评价原始场在全部材料界面与人工拼接界面上的值跳跃和普通法向导数跳跃。正移位预解式的
全界面迹估计把这些跳跃转化为
$\operatorname{dist}(\widehat k^2,\sigma(A(\beta)))$ 的可计算上界；这一全谱距离结论本身
不要求事先验证本质谱隙。只有另行证明所得包含区间落在同一算子的本质谱隙内，才能进一步
把其中的谱点解释为离散导模特征值。本文不据此声称特征值唯一、编号已知或谱隙本身已经
得到证明。第~\ref{sec:modelproblem}~节给出模型，第~\ref{sec:numerics}~节给出候选算法，
第~\ref{sec:error}~节证明谱距离界，第~\ref{sec:results}~节报告一个圆形介质柱、空中心的
谱距离可靠上界的计算实例。
